% from ega3-7-fr.tex
\oldpage[III]{180}
Il est clair que la condition $b')$ de (7.2.5) entraîne $d)$; prouvons que
$d)$ entraîne $b')$.

Il existe un entier $n$ tel que $\mathfrak m^n=0$; posons, pour tout
$A$-module $M$, $M_h=\mathfrak m^hM$; nous prouverons par récurrence
descendante sur $h$ que $t_{M_h}$ est surjectif. La proposition est évidente
pour $h=n$; pour $h<n$, on a une suite exacte
\[
  0\longrightarrow M_{h+1}
  \longrightarrow M_h
  \longrightarrow M_h/M_{h+1}
  \longrightarrow 0
\]
et l'hypothèse de récurrence entraîne que $t_{M_{h+1}}$ est surjectif.
D'autre part, $M_h/M_{h+1}$ est annulé par $\mathfrak m$ et est donc un
$(A/\mathfrak m)$-module, autrement dit est somme directe de $A$-modules
isomorphes à $k$. Pour prouver que $t_{M_h/M_{h+1}}$ est surjectif, il suffit
donc de prouver que $t_k$ l'est, puisque $T$ commute aux sommes directes. Or,
en vertu de la commutativité du diagramme
\[
\xymatrix@C=35pt@R=30pt{
  T(A_s)\otimes_A A_s
  \ar[r]^-{t_{A_s}}\ar[d]_{1\otimes\varepsilon}
  &
  T(A_s)\ar[d]^{T(\varepsilon)}
  \\
  T(A_s)\otimes_A k
  \ar[r]_-{t_k}
  &
  T(k)
}
\]
et de (7.2.4), l'hypothèse $d)$ entraîne que $t_k$ est bien surjectif. Pour
terminer la démonstration, il suffira de prouver que si l'on a une suite
exacte
\[
  0\longrightarrow M'
  \longrightarrow M
  \xrightarrow{v}M''
  \longrightarrow 0
\]
de $A$-modules, telle que $t_{M'}$ et $t_{M''}$ sont surjectifs, alors $t_M$
est surjectif. Or, on a un diagramme commutatif
\[
\xymatrix@C=27pt@R=30pt{
  T'(M')\ar[r]\ar[d]_{t_{M'}}
  &
  T'(M)\ar[r]\ar[d]_{t_M}
  &
  T'(M'')\ar[r]\ar[d]_{t_{M''}}
  &
  0\ar[d]
  \\
  T(M')\ar[r]
  &
  T(M)\ar[r]_-{T(v)}
  &
  T(M'')\ar[r]
  &
  \operatorname{Coker}(T(v))
}
\]
dans lequel les deux lignes sont exactes, en vertu de l'hypothèse que $T$ est
semi-exact. Comme par l'hypothèse de récurrence $t_{M'}$ et $t_{M''}$ sont
des épimorphismes et que la dernière flèche verticale est un monomorphisme,
le lemme des cinq (M, I, 1.1) montre que $t_M$ est un épimorphisme.
\end{proof}

\subsection{Critères d'exactitude des foncteurs homologiques de modules}
\label{subsection:III.7.3.fr}

\begin{proposition}[7.3.1]
Soient $A$ un anneau (non nécessairement commutatif), $T_\bullet$ un foncteur
homologique covariant (T, II, 2.1) de la catégorie $\mathbf{Ab}_A$ dans la
catégorie $\mathbf{Ab}$, commutant aux sommes directes. Soit $p$ un entier tel
que $T_p$ et $T_{p-1}$ soient définis. Les conditions suivantes sont
équivalentes :
\begin{enumerate}
  \item[a)] $T_p$ est exact à droite.
  \item[b)] $T_{p-1}$ est exact à gauche.
  